\subsection{Pairwise Comparison Method}

\content{
        \begin{example} Let's revisit an example from last time. \\
        \begin{tabular}{|l|c|c|c|c|}
        \hline
        & 1 & 3 & 3 & 3 \\
        \hline
        1st Choice  & A & A & O & H \\
        \hline
        2nd Choice  & O & H &H&A \\
        \hline
        3rd Choice & H&O&A&O \\
        \hline
        \end{tabular}
        \end{example}
}{% presentation
\vspace{4in}
}{% nopresentation
\vspace{2in}
}{% comments
Review plurality method. and go over all pairwise comparisons. Is there a
Condorcet Candidate? Did they win the election? Review definitions of Condorcet
Candidate and the Condorcet Criteria if needed. 
}


\content{
        \begin{definition} (Pairwise Comparison Method) Given a preference
        schedule for an election, perform all one-on-one comparisons (eliminate
        all but 2 candidates at a time). A candidate recieves a point for each comparison 
        the candidate wins. If there is a tie, each candidate gets half a point.
        At the end, the candidate with the most points wins. 
        \end{definition}
}{% presentation
}{% nopresentation
}{% comments
}

\content{
        \begin{example} Suppose we must award one of four students, labeled
        $A$, $B$, $C$, and $D$, a scholarship. The comitee submits the following
        preference schedule. With the pairwise comparison method, which student
        will recieve the scholarship? \\
        \begin{tabular}{|l|c|c|c|c|} 
        \hline
        & 5&5&6&4\\
        \hline
        1st & D&A&C&B\\
        \hline
        2nd & A&C&B&D\\
        \hline
        3rd & C&B&D&A\\
        \hline
        4th & B&D&A&C\\
        \hline
        \end{tabular}
        \end{example}
}{% presentation
\vspace{3in}
}{% nopresentation
\newpage
\vspace{1in}
}{% comments
}

\content{
        \begin{example} Your turn! Using the pairwise comparison method, which
        candidate wins the following election? \\
        \begin{tabular}{|l|c|c|c|c|c|c|}
        \hline
        & 3&4&4&6&2&1\\
        \hline
        1st & B&C&B&D&B&E\\
        \hline
        2nd & C&A&D&C&E&A\\
        \hline
        3rd & A&D&C&A&A&D\\
        \hline
        4th & D&B&A&E&C&B\\
        \hline
        5th & E&E&E&B&D&C\\
        \hline
        \end{tabular}\\
        (Hint: there are 10 comparisons to make on this one)
        \end{example}
}{% presentation
\vspace{4in}
}{% nopresentation
\vspace{3in}
}{% comments
}

\content{
        \textbf{Question:} Does the pairwise comparison method satisfy the
        Condorcet Criterion?
}{% presentation
\vspace{1in}
}{% nopresentation
\vspace{2cm}
}{% comments
}

\content{
        \begin{example} Using the pairwise comparison method, which
        candidate wins the following election? \\
        \begin{tabular}{|l|c|c|c|c|}
        \hline
        & 13&14&11&9\\
        \hline
        1st & B&D&B&C\\
        \hline
        2nd & C&C&A&A\\
        \hline
        3rd & A&A&D&D\\
        \hline
        4th & D&B&C&B\\
        \hline
        \end{tabular}\\
        \end{example}
}{% presentation
\vspace{1in}
}{% nopresentation
\vspace{2in}
\newpage
}{% comments
}


\content{
        So what is wrong with this method? 
        \begin{example} Suppose in one of the previous examples, (recall that
        $C$ was the winner), the candidate $D$ is discovered to be inelegible
        for the scholarship due to failing a class last semester. With $D$ gone,
        the comittee decides to recount the votes using the same method. Is the
        outcome the same? \\
        \begin{tabular}{|l|c|c|c|c|} 
        \hline
        & 5&5&6&4\\
        \hline
        1st & D&A&C&B\\
        \hline
        2nd & A&C&B&D\\
        \hline
        3rd & C&B&D&A\\
        \hline
        4th & B&D&A&C\\
        \hline
        \end{tabular}
        \end{example}
}{% presentation
\vspace{4in}
}{% nopresentation
\vspace{2in}
}{% comments
}


\content{
        \begin{definition} (Independence of Irrelevant Alternatives (IIA)
        Criterion) If a non-winning candidate is removed from the ballot, it
        should not change the winner of the election. 
        \end{definition}
}{% presentation
\vspace{1in}
}{% nopresentation
}{% comments
Note that the pairwise comparison method can violate this criterion. 
}
